\chapter{Boot Sequence}

\section{Overview}

QSOE boots through a well-defined chain of handoffs, from firmware
through the elfloader into the seL4 kernel, which then starts the
initial user-space process (the \emph{root server}).

\begin{figure}[ht]
\centering
\begin{tikzpicture}[
    stage/.style={
        draw, rounded corners=3pt, minimum width=10cm,
        minimum height=1.0cm, align=center, font=\small
    },
    arr/.style={-{Stealth[length=6pt]}, thick},
]
    \node[stage, fill=gray!10]                          (sbi)  {OpenSBI (M-mode firmware)};
    \node[stage, fill=green!8,  below=0.5cm of sbi]    (elf)  {elfloader --- extracts kernel.elf + rootserver.elf from CPIO};
    \node[stage, fill=blue!8,   below=0.5cm of elf]    (kern) {seL4 kernel --- creates initial thread, sets PC to rootserver entry};
    \node[stage, fill=orange!8, below=0.5cm of kern]   (rt)   {\texttt{\_sel4\_start} (sel4runtime) $\rightarrow$ \texttt{main()}};

    \draw[arr] (sbi)  -- node[right, font=\footnotesize] {a0=hart\_id, a1=DTB} (elf);
    \draw[arr] (elf)  -- node[right, font=\footnotesize] {a0\ldots a7: rootserver info} (kern);
    \draw[arr] (kern) -- node[right, font=\footnotesize] {\texttt{sret} to user space} (rt);
\end{tikzpicture}
\caption{QSOE boot sequence on RISC-V}
\label{fig:boot}
\end{figure}

\section{OpenSBI}

OpenSBI runs in M-mode and provides the Supervisor Binary Interface.
It initializes hardware, handles traps that must be served in M-mode,
and jumps to the elfloader at the address specified in the firmware
configuration.  It passes the hart ID in \texttt{a0} and the device
tree blob (DTB) address in \texttt{a1}.

\section{Elfloader (startup)}

The elfloader is a minimal S-mode program whose sole purpose is to
unpack and launch the kernel.  It contains an embedded CPIO archive
with two ELF images:

\begin{itemize}
    \item \texttt{kernel.elf} --- the seL4 microkernel
    \item \texttt{rootserver.elf} --- the initial user-space process
          (taskman, linked with sel4runtime and musl libc)
\end{itemize}

The elfloader:
\begin{enumerate}
    \item Parses the CPIO archive to locate both ELF images.
    \item Loads each image into its designated physical memory region.
    \item Reads the ELF entry points (\texttt{e\_entry}) from both images.
    \item Sets up initial page tables for the kernel.
    \item Jumps to the kernel entry point, passing the rootserver's
          physical location and virtual entry address via registers
          \texttt{a0}--\texttt{a7}.
\end{enumerate}

\section{seL4 kernel initialization}

The kernel receives the rootserver information through the standard
RISC-V calling convention:

\begin{table}[ht]
\centering
\begin{tabular}{@{}ll@{}}
    \toprule
    \textbf{Register} & \textbf{Content} \\
    \midrule
    \texttt{a0} & rootserver physical start \\
    \texttt{a1} & rootserver physical end \\
    \texttt{a2} & physical--virtual offset \\
    \texttt{a3} & rootserver virtual entry (\texttt{\_sel4\_start}) \\
    \texttt{a4} & DTB address \\
    \texttt{a5} & DTB size \\
    \texttt{a6} & hart ID \\
    \texttt{a7} & core ID \\
    \bottomrule
\end{tabular}
\caption{Registers passed from elfloader to kernel}
\label{tab:bootregs}
\end{table}

The kernel creates the initial thread (root task), stores the
rootserver's entry address via \texttt{setNextPC()}, populates the
\texttt{seL4\_BootInfo} structure with available capabilities and
untyped memory, and executes \texttt{sret} to enter user space.

\section{Root server startup (sel4runtime)}

Execution lands at \texttt{\_sel4\_start} in
\texttt{crt/arch/riscv/sel4\_crt0.S}, which:
\begin{enumerate}
    \item Sets the global pointer (\texttt{gp}) for linker relaxation.
    \item Establishes a 16\,KB stack.
    \item Calls \texttt{\_\_sel4\_start\_root(boot\_info)} in C, which
          parses the auxiliary vector, initializes TLS via the \texttt{tp}
          register, runs constructors, and finally calls \texttt{main()}.
\end{enumerate}

At this point, \texttt{taskman}'s \texttt{main()} takes over.
